% =========================
% 3_motivation.tex
% =========================
\section{Motivation}
\label{sec:motivation}

Section~\ref{sec:why} argues from concrete deployments why open-set labels, sensing
heterogeneity, and a small supervision budget are the binding constraints, rather than
features a paper might optionally offer. Section~\ref{sec:heterogeneity_motiv} then
characterises the heterogeneity technically. Section~\ref{sec:no_grammar} argues that it
cannot be closed by data---inertial signals admit no closed grammar---so compatibility must
be supplied as inductive bias, and that the bias must be \emph{calibrated} rather than
maximised. Section~\ref{sec:geometry} shows by direct measurement that the field's standard
mechanism for open-ended labels rests on a false premise. Section~\ref{sec:requirements}
turns these into design requirements.

\subsection{Why This Problem: Three Deployments and What They Demand}
\label{sec:why}

That IMU-based activity recognition is useful needs no argument. What needs an argument is why
it should be \emph{open-set}, \emph{heterogeneity-aware}, and \emph{label-efficient}---three
qualifiers that add real difficulty and that a system could reasonably decline to take on. Our
answer is that in the deployments that motivate the work, all three are forced simultaneously,
and dropping any one of them removes the application rather than simplifying it.

\textbf{Rehabilitation adherence.} A physiotherapist prescribes a set of exercises after
surgery and needs to know whether they were performed, and performed correctly. The exercise
set is chosen \emph{per patient} from a long clinical tail---a specific loaded knee extension,
a specific scapular retraction---so no fixed output vocabulary covers it. The sensor is the
patient's own phone or watch, worn wherever they choose. The available supervision is whatever
the clinician can demonstrate inside a twenty-minute appointment.

\textbf{Occupational strain.} A safety team wants to count a specific risky motion in one
facility: reaching under a chassis, lifting from below knee height. The vocabulary belongs to
that facility and its risk assessment. The hardware is a fleet accumulated over years, mixing
phone models and wristbands. Collecting a labelled corpus would require instrumenting the floor
with video and obtaining consent from every worker in frame---which is precisely the thing the
deployment exists to avoid.

\textbf{Free-living behavioural monitoring.} A cohort study measures posture, sedentary time,
and sleep over weeks. Devices vary across participants and firmware revisions vary within a
participant; the categories are the study's own, defined by its protocol rather than by a
benchmark.

\begin{table}[t]
  \centering
  \resizebox{\columnwidth}{!}{
  \small
  \setlength{\tabcolsep}{4pt}
  \begin{tabular}{l ccc}
    \toprule
    & open-set & heterogeneous & small label \\
    Deployment & labels & sensing & budget \\
    \midrule
    Rehabilitation adherence   & \checkmark & \checkmark & \checkmark \\
    Occupational strain        & \checkmark & \checkmark & \checkmark \\
    Free-living monitoring     & \checkmark & \checkmark & \checkmark \\
    \midrule
    \emph{Public HAR benchmark} & --- & --- & --- \\
    \bottomrule
  \end{tabular}}
  \caption{The three qualifiers are not independent features to be offered separately. Every
  deployment that motivates this work demands all three at once, because all three are
  properties of the deployment rather than of the recognition task. A public benchmark demands
  none of them---which is why a model can score well on one and still be unusable in any of the
  rows above.}
  \label{tab:demands}
  \vspace{-1em}
\end{table}

Table~\ref{tab:demands} makes the pattern explicit, and also makes clear why benchmark
performance has been a poor guide: a fixed-vocabulary, fixed-layout, fully-supervised benchmark
demands none of the three, so a model can lead it and still be inapplicable to every row above.
We now take each qualifier in turn, because each rests on a different argument.

\textbf{Why open-set is not optional.} Two reasons, neither of which is solved by more data.
First, \emph{the label set is not the model's to choose}. A prescribed rehabilitation exercise
is defined per patient; a task vocabulary is defined per facility. These are not rare classes
that a larger corpus would eventually cover---they are defined after training, by someone who
is not us. A model with a fixed output layer is structurally unable to express them.
Second, \emph{the annotation channel for inertial data is uniquely poor}. Images and text can
be labelled from the artifact itself, which is why annotation crowdsources and why those
corpora are enormous. Nobody can label an accelerometer trace by looking at it. Labels must
therefore come from synchronised video, scripted in-laboratory protocols, or self-report---all
expensive, and all biasing the resulting data toward laboratory-like conditions. This is the
mechanical reason HAR corpora are small and scripted while image corpora are large and natural,
and it will not be fixed by effort. A system that requires a labelled target set per deployment
inherits that cost every time.

\textbf{Why heterogeneity-aware is not optional.} The cost of sensing heterogeneity is usually
described as a one-time integration expense. It is not: it \emph{recurs}, per device model, per
placement, per firmware revision, for as long as the product ships. A vendor covering fifty
device-by-placement combinations either maintains fifty models---fifty training runs, fifty
validation cycles, fifty entries in a model registry---or ships one model built for the lowest
common denominator. In practice the second happens, and it is a large part of why deployed
phone activity recognition is still limited to a handful of coarse ambulatory classes. It is
also why there is no ``download the model and point it at your sensor'' moment for IMU as there
is for vision: the model cannot read your sensor. Notably, this is not only our reading.
UniMTS~\cite{unimts}, the leading zero-shot motion model, ablates its own components and
attributes substantially more of its performance to acquisition-configuration handling---
rotation-invariant augmentation and placement encoding---than to its label-text handling. The
field names sensor heterogeneity and poor activity text as its barriers; it does not name label
semantics.

\textbf{Why the label budget is the right axis.} Given the two arguments above, the useful
question is not whether a system is ``zero-shot'' but how much target supervision it needs to
become useful, and what that supervision costs to obtain. Zero labelled examples is the
interesting left endpoint because it is the only point reachable without any target data
collection at all. But a clinician demonstrating an exercise five times is a realistic and cheap
form of supervision, and a system that can consume it \emph{without a training cycle} is
categorically easier to deploy than one that requires fine-tuning: nothing needs to be
retrained, nothing needs re-validating, and the demonstration never has to leave the device.
That last property is a legal precondition in the clinical and workplace settings above, not a
convenience.

\textbf{What this buys, and why it should interest a systems audience.} The result of meeting
all three is a difference in \emph{capability}, not a difference in accuracy: one encoder
serving many deployments instead of one model per configuration; new activities and new devices
added by appending to a memory rather than by retraining, so that a model already validated does
not have to be validated again; enrollment performed on-device; and predictions that can be
inspected, because a retrieval-based decision can name the stored examples that produced it. In
a regulated or clinical setting the last two are worth more than a few points of macro-F1.

\subsection{Sensing Heterogeneity Across Devices and Deployments}
\label{sec:heterogeneity_motiv}

Wearable and mobile devices differ widely in their sensing configurations, and the
differences are not cosmetic. Our training corpus alone (Table~\ref{tab:datasets},
Figure~\ref{fig:corpus}) spans sampling rates from $20$ to $100$\,Hz, accelerometer-only and
accelerometer-plus-gyroscope suites, and sensor placements at the waist, trouser pocket,
wrist, lower back, ear, and head---the last two from smart glasses and wireless earbuds,
devices that did not exist as HAR platforms when most public benchmarks were collected.

\textbf{Sampling rate.} At $20$\,Hz a $6$-second window is $120$ samples; at $100$\,Hz it is
$600$. Any encoder whose positional code is tied to a timestep index conflates physical time
with sequence length, so the same motion at two rates presents as two different structures.
The usual fix---resample everything to a common rate---is worse than it looks: downsampling
to the corpus minimum aliases the high-frequency content that distinguishes, say, brushing
teeth from waving, and upsampling fabricates samples that were never measured.

\textbf{Channel set, placement, and orientation.} Configurations range from a single 3-axis
wrist accelerometer to six simultaneous placements on one body. A fixed-width input layer
cannot accommodate this without padding artifacts or per-dataset input heads. Worse, the
\emph{meaning} of a channel depends on where and how the device is worn: a wrist
accelerometer's $x$-axis and a hip accelerometer's $x$-axis are not the same physical
quantity, and neither is stable if the user pockets the phone upside down.

\textbf{Conventions the data does not announce.} Two accelerometer streams that look
identical may differ in gravity convention (Android reports total acceleration; iOS
\texttt{userAcceleration} removes gravity), in units, or in an unadvertised on-device
low-pass filter. In assembling our corpus we found several public datasets whose distributed
form had already been transformed---one is per-axis min--max normalised, destroying physical
scale; another ships gravity-removed body acceleration under a name suggesting total
acceleration. A model that silently assumes a convention will be wrong on data it has every
reason to think it understands.

\begin{figure}[t]
    \centering
    \setlength{\abovecaptionskip}{2pt}
    \setlength{\belowcaptionskip}{-4pt}
  % =====================================================================
% fig_corpus.tex  --  corpus heterogeneity: placement x sampling rate
%
% DATA PROVENANCE: every point is a real stream in the materialized native
% grids, enumerated 2026-07-26 from training.tokenizer.pretrain_data.CorpusIndex
% (20 training streams over 12 datasets) and data/datasets/*/metadata.json
% (7 held-out datasets).  Bubble area encodes window count; the largest
% (CAPTURE-24 wrist) is 144,142 windows, the smallest (SP-SW-HAR pocket) 1,179.
% =====================================================================
\resizebox{\columnwidth}{!}{%
\begin{tikzpicture}
\begin{axis}[
  width=\columnwidth, height=4.6cm,
  xmin=0.55, xmax=5.55,
  ymin=6, ymax=145,
  ytick={20,50,100},
  yticklabels={20\,Hz,50\,Hz,100\,Hz},
  yticklabel style={font=\scriptsize},
  xtick={1,2,3,4,5},
  xticklabels={waist/hip,pocket,wrist,lower back,head/ear},
  xticklabel style={font=\scriptsize},
  ylabel={sampling rate},
  ylabel style={font=\scriptsize,yshift=-6pt},
  grid=major, grid style={draw=halogray!18},
  axis line style={draw=halogray!60},
  tick style={draw=halogray!60},
  scale only axis,
  clip=false,
]
% ---------- training streams (filled blue discs, diameter ~ #windows^0.4) ----
% diameters below are 2*(1.2 + 4.5*(n/144142)^0.4) pt, n = window count
\tikzset{hdisc/.style={circle,draw=haloblue,fill=haloblue!45,line width=0.4pt,inner sep=0pt}}
\node[hdisc,minimum size=5.44pt] at (axis cs:0.88,50)  {};% hhar        9,587
\node[hdisc,minimum size=5.40pt] at (axis cs:1.00,100) {};% kuhar       9,233
\node[hdisc,minimum size=5.54pt] at (axis cs:1.12,50)  {};% uci_har    10,299
\node[hdisc,minimum size=5.70pt] at (axis cs:1.80,50)  {};% unimib     11,771
\node[hdisc,minimum size=4.88pt] at (axis cs:1.94,50)  {};% xrf l.pkt   5,794
\node[hdisc,minimum size=7.76pt] at (axis cs:2.00,20)  {};% wisdm pkt  39,329
\node[hdisc,minimum size=4.88pt] at (axis cs:2.06,50)  {};% xrf r.pkt   5,794
\node[hdisc,minimum size=3.72pt] at (axis cs:2.00,100) {};% spsw pkt    1,179
\node[hdisc,minimum size=11.4pt] at (axis cs:2.82,100) {};% capture24 144,142
\node[hdisc,minimum size=6.36pt] at (axis cs:2.85,50)  {};% harmes     18,576
\node[hdisc,minimum size=4.36pt] at (axis cs:2.94,100) {};% pamap2      3,186
\node[hdisc,minimum size=3.68pt] at (axis cs:2.95,50)  {};% mhealth     1,104
\node[hdisc,minimum size=7.26pt] at (axis cs:3.00,20)  {};% wisdm wr   30,876
\node[hdisc,minimum size=4.88pt] at (axis cs:3.05,50)  {};% xrf l.wr    5,794
\node[hdisc,minimum size=5.86pt] at (axis cs:3.06,100) {};% nfi wrist  13,260
\node[hdisc,minimum size=4.88pt] at (axis cs:3.15,50)  {};% xrf r.wr    5,794
\node[hdisc,minimum size=4.20pt] at (axis cs:3.18,100) {};% spsw wrist  2,578
\node[hdisc,minimum size=5.86pt] at (axis cs:4.00,100) {};% nfi back   13,260
\node[hdisc,minimum size=4.88pt] at (axis cs:4.92,50)  {};% xrf glasses 5,794
\node[hdisc,minimum size=4.88pt] at (axis cs:5.08,50)  {};% xrf earbuds 5,794
% ---------- held-out datasets (open orange triangles) ----------
\addplot[only marks,mark=triangle*,mark size=2.6pt,color=haloorange,
         fill=halolight2,line width=0.7pt]
  coordinates {
    (1.26,50) (1.38,50) (1.26,100)
    (2.20,50) (2.32,50)
    (3.30,50) (3.42,50)
  };
% ---------- annotations ----------
\node[font=\scriptsize,color=haloblue!75!black,anchor=west] at (axis cs:2.95,128)
  {\textbf{CAPTURE-24} 144k windows};
\draw[haloblue!70,thin] (axis cs:2.86,106) -- (axis cs:2.93,124);
\node[font=\scriptsize,color=halogray,anchor=west,align=left] at (axis cs:4.05,24)
  {smart glasses\\ \& earbuds};
\draw[halogray,thin] (axis cs:5.00,44) -- (axis cs:4.90,31);
\node[font=\scriptsize,anchor=west,align=left] at (axis cs:0.62,13)
  {\textcolor{haloblue!75!black}{$\bullet$ 20 training streams}\;\;
   \textcolor{haloorange!85!black}{$\blacktriangle$ 7 held-out}};
\end{axis}
\end{tikzpicture}}
\caption{\textbf{The corpus does not cover a grid; it covers a scatter.} Each blue disc is
one of the 20 sensor streams HALO trains on, positioned by body placement and native
sampling rate, with area proportional to the number of $6$\,s windows it contributes; orange
triangles are the 7 held-out evaluation datasets. Rates span a $5\times$ range and placements
include two---smart glasses and wireless earbuds---that no established HAR benchmark
contains. Held-out cells are not simply new subjects: they combine placements and rates in
proportions the training corpus never sees. Absorbing this by resampling to a common rate
would discard the content above $10$\,Hz that only the $50$/$100$\,Hz streams can
resolve.}
\label{fig:corpus}

\end{figure}

\subsection{Inertial Signals Admit No Closed Grammar}
\label{sec:no_grammar}

The heterogeneity above is usually treated as an engineering nuisance to be absorbed by
pretraining on more of it. We think that is the wrong diagnosis, and the difference matters
because it changes what the solution has to be.

Foundation models in language and vision work in large part because their input space is
effectively \emph{closed} by a corpus. Text is produced by a grammar over a finite lexicon;
enough text does approach the distribution of sentences a model will meet. Natural images
obey a stable projective geometry and a stable statistics of surfaces and light. In both
cases the substrate has structure that a sufficiently large sample covers.

Inertial signals have no such closure. The map from a motion to a recorded signal is
composed at deployment time from a chain of choices---mounting orientation (a continuous
$SO(3)$ nuisance), body placement, sampling rate, filter chain, gravity convention, sensor
suite---and the chain composes. There is no finite enumeration to cover, and the free
parameters are continuous. A corpus $10\times$ larger reduces the \emph{probability} that a
new deployment is far from anything seen, but the tail does not close, and the transformations
that generate it are cheap to encounter: shipping on a second phone model, or a firmware
update, is enough.

If compatibility cannot come from coverage, it must come from \emph{inductive bias}: the
architecture, not the data, must make the representation invariant to the nuisance axes and
explicitly aware of the ones that carry signal. This is the standard move---but for inertial
sensing it has a failure mode on \emph{both} sides, and the interesting design question is
where to sit between them.

\textbf{Too much bias.} Layout-locked encoders fix a channel order, a channel count, and a
rate. They are the most sample-efficient thing to train and cannot read a device they were
not built for. Orientation-\emph{invariant} designs (e.g.\ training against uniformly
sampled $SO(3)$ rotations) are a subtler version of the same error: invariance is lossy.
Discarding orientation discards gravity, and with it the distinction between lying down and
standing still, which is exactly the kind of activity a posture or sleep application cares
about.

\textbf{Too little bias.} The opposite extreme treats every scalar channel as an independent
univariate series identified only by a free-text name. This is maximally flexible, and it
creates a shortcut. In a multi-corpus training set, the string describing a channel is
nearly a sufficient statistic for \emph{which corpus} the window came from; so is the
sampling rate, and so is the channel count. A model rewarded for predicting a label within a
corpus can satisfy much of that objective by first inferring the corpus and then predicting
its marginal---and cross-dataset evaluation is precisely where that strategy stops working.
We therefore treat corpus identifiability as a measurable failure and probe for it directly
(Table~\ref{tab:identity_probe}).

A third position, recently explored outside HAR, is to make the model \emph{equivariant} to
resolution---to condition its computation on the sampling rate so that it adapts, as in
sampling-rate-equivariant forecasting~\cite{flowstate} or grid-agnostic spectral
tokenization~\cite{spectraltok}. That is strictly better than resampling, but it is not the same
as what we want. An equivariant model still has to be \emph{told} the rate, and its features at
two rates are related by a learned transformation rather than being the same quantity. We want
the stronger property: that a feature \emph{denotes} a physical frequency, so a band computed
from a $20$\,Hz phone and a band computed from a $100$\,Hz research IMU are directly comparable
without any conditioning at all---which is what makes them comparable in a shared memory. Within
IMU sensing this property is, as far as we know, unclaimed: the closest tokenizer,
MoPFormer~\cite{mopformer}, resamples every dataset to $100$\,Hz before quantising into motion
primitives, so its codebook is defined over a fixed-rate grid.

The calibrated position, which HALO adopts, is to commit to the physics that does not vary
and to nothing that does. Three things do not vary: co-located axes on one rigid device
share a coordinate frame and must be rotated together, not independently; human motion
energy lives at physical frequencies, so a band means the same thing at every sampling rate;
and the quasi-static component of an accelerometer is gravity, so it carries posture and must
be preserved as a signed quantity rather than normalised away. Everything else---rate,
channel count, ordering, placement, which sensors exist---is deferred to run time and
described in language, at the granularity of the \emph{sensor} rather than the channel, so
that the description is expressive enough to name an unseen device but too coarse to
fingerprint a corpus.

\subsection{Sensor Similarity Does Not Track Linguistic Similarity}
\label{sec:geometry}

The second obstacle concerns the output side. Because deployments name their own activities,
a reusable model cannot have a fixed output layer, and the field's answer has been to embed
label strings with a frozen sentence encoder and align sensor features to that space, so
that recognition becomes nearest-neighbour retrieval over an arbitrary label
library~\cite{clip,lanhar,goat,unimts}. The mechanism is elegant and its premise is rarely
stated: that proximity in sentence-embedding space is a usable proxy for proximity in motion
space.

% =====================================================================
% fig_geometry.tex  --  L2 evidence: sensor similarity != linguistic similarity
%
% DATA PROVENANCE: cosine similarities between all-MiniLM-L6-v2 sentence
% embeddings of the label strings, measured 2026-07-26 by
%   sentence_transformers.SentenceTransformer('all-MiniLM-L6-v2')
% with normalize_embeddings=True.  Raw values in the arrays below; the
% generating script and its JSON output are archived with the paper source.
% The MOTION grouping (same vs different physical motion) is our own
% judgement and is stated as such in the caption -- it is not measured.
% =====================================================================
\begin{figure}[t]
\centering
\setlength{\abovecaptionskip}{3pt}
\setlength{\belowcaptionskip}{-4pt}
\resizebox{\columnwidth}{!}{%
\begin{tikzpicture}
\begin{axis}[
  width=\columnwidth, height=4.35cm,
  xmin=0.40, xmax=1.02,
  ymin=-1.15, ymax=1.85,
  xtick={0.4,0.5,0.6,0.7,0.8,0.9,1.0},
  xticklabel style={font=\scriptsize},
  xlabel={cosine similarity of the two label strings (frozen sentence encoder)},
  xlabel style={font=\scriptsize,yshift=2pt},
  ytick={0,1},
  yticklabels={},
  axis y line=none,
  axis x line*=bottom,
  clip=false,
  scale only axis,
]
% ---- shaded overlap region: where the two groups interleave ----
\addplot[draw=none,fill=halogray!12,forget plot]
  coordinates {(0.5239,-1.15) (0.7965,-1.15) (0.7965,1.85) (0.5239,1.85)}
  \closedcycle;
\node[font=\scriptsize\itshape,halogray,anchor=south,align=center]
  at (axis cs:0.660,1.50) {overlap};

% ---- lane guides ----
\addplot[draw=halogray!35,thin,forget plot] coordinates {(0.40,1) (1.02,1)};
\addplot[draw=halogray!35,thin,forget plot] coordinates {(0.40,0) (1.02,0)};

% ---- TOP LANE: DIFFERENT physical motion, near-identical text ----
\addplot[only marks,mark=*,mark size=1.9pt,color=haloorange,forget plot]
  coordinates {(0.5239,1) (0.6396,1) (0.7668,1) (0.7690,1) (0.8678,1) (0.8990,1) (0.9342,1)};
% group mean
\addplot[only marks,mark=|,mark size=5pt,line width=1.1pt,color=haloorange!70!black,forget plot]
  coordinates {(0.7715,1)};

% ---- BOTTOM LANE: SAME physical motion, distant text ----
\addplot[only marks,mark=*,mark size=1.9pt,color=haloblue,forget plot]
  coordinates {(0.4500,0) (0.4595,0) (0.4917,0) (0.5599,0) (0.6645,0) (0.7775,0) (0.7965,0)};
\addplot[only marks,mark=|,mark size=5pt,line width=1.1pt,color=haloblue!70!black,forget plot]
  coordinates {(0.5999,0)};

% ---- lane titles ----
\node[font=\scriptsize\bfseries,color=haloorange!75!black,anchor=west]
  at (axis cs:0.405,1.42) {different motion, near-identical text};
\node[font=\scriptsize\bfseries,color=haloblue!75!black,anchor=west]
  at (axis cs:0.405,-0.72) {same motion, distant text};

% ---- callouts on the two headline pairs ----
\node[font=\scriptsize,anchor=south east,inner sep=1pt]
  (sitdown) at (axis cs:1.015,1.06) {\emph{sitting} / \emph{sitting down}};
\node[font=\scriptsize,anchor=north west,inner sep=1pt]
  (coffee) at (axis cs:0.500,-0.20) {\emph{drinking coffee} / \emph{drinking water}};
\draw[haloblue,thin] (axis cs:0.5599,-0.10) -- (axis cs:0.552,-0.30);

% ---- mean labels ----
\node[font=\scriptsize,color=haloorange!70!black,anchor=south] at (axis cs:0.7715,1.10) {$\bar{s}=.77$};
\node[font=\scriptsize,color=haloblue!70!black,anchor=south]   at (axis cs:0.5999,0.10) {$\bar{s}=.60$};
\end{axis}
\end{tikzpicture}}
\caption{\textbf{Sensor similarity and linguistic similarity are inverted, not merely
noisy.} Each dot is one pair of activity names from our label vocabulary, placed at the
cosine similarity of their frozen sentence embeddings (all-MiniLM-L6-v2, the encoder
language-aligned HAR models use). Pairs in the \emph{top} lane denote physically distinct
activities; pairs in the \emph{bottom} lane denote the same physical motion. If text
similarity tracked sensor similarity the lanes would separate, with the bottom lane on the
right. They separate the wrong way: the mean of the physically \emph{distinct} pairs
($0.77$) exceeds the mean of the physically \emph{identical} pairs ($0.60$), and $38$ of the
$49$ cross-lane comparisons are inverted. The motion grouping is our judgement; the
similarities are measured. A single projection from sensor space into this text geometry is
therefore mis-specified---which is why HALO retrieves in sensor space and votes in label
space (Section~\ref{sec:evidence_design}).}
\label{fig:geometry}
\end{figure}


It is not. Figure~\ref{fig:geometry} measures this on our own vocabulary. We selected $14$
label pairs---seven where the two names denote the same physical motion, seven where they
denote different motions---and computed the cosine similarity of their embeddings under
all-MiniLM-L6-v2, the encoder these systems use. If the premise held, the same-motion pairs
would sit to the right of the different-motion pairs. The ordering is reversed: the
different-motion pairs average $0.77$ and the same-motion pairs $0.60$, and $38$ of the $49$
cross-comparisons are individually inverted. \emph{Drinking coffee} and \emph{drinking
water} are the same wrist trajectory at $0.56$; \emph{sitting} and \emph{sitting down}, a
static posture versus a dynamic transition, are at $0.93$.

The failure is structural rather than a deficiency of one sentence encoder. Natural language
encodes activities by their \emph{object and purpose}---what is being drunk, what is being
cut---while an IMU sees only the \emph{kinematics of the limb}, which is invariant to the
object. Conversely, morphological relatives (\emph{sit}/\emph{sit down},
\emph{upstairs}/\emph{downstairs}) share nearly all of their surface form while denoting
opposite dynamics. No amount of alignment training repairs a target geometry that is wrong;
a model forced to fit it will minimise its loss by other means, and within a multi-dataset
corpus the cheapest other means is dataset identity.

The architectural consequence is to stop trying. HALO keeps language as the \emph{interface}
for naming activities---that part is necessary, and cheap---but does not use it as the
\emph{metric space} in which recognition happens. Retrieval is performed in sensor space,
where similarity means kinematic similarity; label text enters only afterwards, to name the
retrieved neighbours and to enumerate what the deployment is willing to predict
(Section~\ref{sec:evidence_design}).

\subsection{Design Requirements}
\label{sec:requirements}

The three claims yield four requirements.

\textbf{R1 (from L1).} The encoder must accept arbitrary sampling rates, channel counts,
sensor suites, and placements with one weight set and no architectural change, absorbing them
through construction rather than through resampling or per-dataset heads---and must signal
what a given configuration could not observe instead of silently fabricating it.

\textbf{R2 (from L1).} The conditioning that carries acquisition context must generalise
compositionally to unseen configurations, and must be coarse enough that it cannot serve as
a corpus identifier. This is a testable property, not an assertion.

\textbf{R3 (from L2).} Recognition over a run-time-declared label set must not require a
single projection reconciling sensor geometry with sentence geometry. Similarity must be
computed where it is meaningful, and label text used only where it is meaningful.

\textbf{R4 (from L3).} The model must place no floor on how much target supervision is
required, and adding a small amount must be possible without a gradient update, without
retraining, and without data leaving the device---so that $k{=}0$, $k{=}10$, and full
supervision are points on one curve rather than three different systems.

HALO meets R1--R2 with the physical-frequency filterbank tokenizer and factored per-sensor
language conditioning (Section~\ref{sec:tokenizer_design}), R3 with the evidence engine
(Section~\ref{sec:evidence_design}), and R4 because the evidence engine's memory is
append-only (Section~\ref{sec:deployment}).
